% !TEX root = ../thesis.tex
% ==============================================================================
%  PERSIAN ABSTRACT & KEYWORDS (frontmatter/abstract.tex)
% ==============================================================================
%
%  This file contains the Persian abstract and keywords (چکیده و کلیدواژه‌ها).
%
% ------------------------------------------------------------------------------
%  HOW TO USE & CUSTOMIZE:
% ------------------------------------------------------------------------------
%  - Wrap your abstract text in \begin{persianabstract} ... \end{persianabstract}.
%  - Add your Persian keywords using \thesiskeywords{کلیدواژه اول، کلیدواژه دوم، ...}.
%  - Vertical spacing above the abstract can be adjusted in `styles/custom.tex`
%    using \setlength{\thesisabstractbeforespace}{...}.
%  - To toggle this page in the thesis:
%    Toggle `% !TEX root = ../thesis.tex
% ==============================================================================
%  PERSIAN ABSTRACT & KEYWORDS (frontmatter/abstract.tex)
% ==============================================================================
%
%  This file contains the Persian abstract and keywords (چکیده و کلیدواژه‌ها).
%
% ------------------------------------------------------------------------------
%  HOW TO USE & CUSTOMIZE:
% ------------------------------------------------------------------------------
%  - Wrap your abstract text in \begin{persianabstract} ... \end{persianabstract}.
%  - Add your Persian keywords using \thesiskeywords{کلیدواژه اول، کلیدواژه دوم، ...}.
%  - Vertical spacing above the abstract can be adjusted in `styles/custom.tex`
%    using \setlength{\thesisabstractbeforespace}{...}.
%  - To toggle this page in the thesis:
%    Toggle `% !TEX root = ../thesis.tex
% ==============================================================================
%  PERSIAN ABSTRACT & KEYWORDS (frontmatter/abstract.tex)
% ==============================================================================
%
%  This file contains the Persian abstract and keywords (چکیده و کلیدواژه‌ها).
%
% ------------------------------------------------------------------------------
%  HOW TO USE & CUSTOMIZE:
% ------------------------------------------------------------------------------
%  - Wrap your abstract text in \begin{persianabstract} ... \end{persianabstract}.
%  - Add your Persian keywords using \thesiskeywords{کلیدواژه اول، کلیدواژه دوم، ...}.
%  - Vertical spacing above the abstract can be adjusted in `styles/custom.tex`
%    using \setlength{\thesisabstractbeforespace}{...}.
%  - To toggle this page in the thesis:
%    Toggle `% !TEX root = ../thesis.tex
% ==============================================================================
%  PERSIAN ABSTRACT & KEYWORDS (frontmatter/abstract.tex)
% ==============================================================================
%
%  This file contains the Persian abstract and keywords (چکیده و کلیدواژه‌ها).
%
% ------------------------------------------------------------------------------
%  HOW TO USE & CUSTOMIZE:
% ------------------------------------------------------------------------------
%  - Wrap your abstract text in \begin{persianabstract} ... \end{persianabstract}.
%  - Add your Persian keywords using \thesiskeywords{کلیدواژه اول، کلیدواژه دوم، ...}.
%  - Vertical spacing above the abstract can be adjusted in `styles/custom.tex`
%    using \setlength{\thesisabstractbeforespace}{...}.
%  - To toggle this page in the thesis:
%    Toggle `\include{frontmatter/abstract}` in `thesis.tex`.
% ==============================================================================

\begin{persianabstract}
به‌کارگیری مدل‌های زبانی بزرگ از طریق واسط‌های برنامه‌نویسی مسدود در حوزه‌های پرمخاطره، تفسیرپذیری را ضروری ساخته است. با این حال، پارادایم‌های متداول نارسایی‌های اساسی دارند: رویکردهای درونی با تحمیل قیدهای ساختاری، ظرفیت مدل را محدود کرده و برای سیستم‌های منجمد ناکارآمدند؛ روش‌های پساعملی جعبه‌سفید به اطلاعات درونی پنهان در پسِ واسط‌های تجاری وابسته‌اند؛ و روش‌های پساعملی جعبه‌سیاه نیز یا با استدلال‌های هدفمند در خودتوضیحی مواجه‌اند یا در رویکردهای اختلالی توکن‌محور، به فروپاشی دستوری، خطای خارج از توزیع، عدم درک واژگان هم‌معنی و هزینه‌های نمایی دچار می‌شوند. افزون بر این، پیچیدگی ماهوی مدل‌های مولد سبب شده است تا انتساب ویژگی در رژیم کاملاً جعبه‌سیاه و تنها بر مبنای متن خام، همچنان مسئله‌ای بررسی‌نشده در آن‌ها باقی بماند.

برای حل این مسائل، این پایان‌نامه واحد تحلیل را از توکن‌ها به جملات ارتقا می‌دهد تا با حفظ یکپارچگی دستوری در حذف ویژگی‌ها، توضیحاتی خوانا تولید شوند. افزون بر این، یک مدل مبتنی بر انرژی به عنوان جانشین، هماهنگی پرامپت و پاسخ را بدون نیاز به توابع محاسبه‌ناپذیر می‌سنجد. برای آموزش این چشم‌انداز انرژی در غیاب برچسب‌های حقیقت پایه، راهبرد نمونه‌برداری تقابلی دوگانه‌ای طراحی شده است که وفاداری سراسری (تقابل با متون انسانی و مدل‌های خارجی) را با وابستگی‌های محلی پادواقعی (ماسک‌گذاری ساختاریافته‌ی پرامپت و پاسخ) ترکیب می‌کند.

بر این مبنا، چارچوب \lr{ESCI} به عنوان روشی پساعملی و مبتنی بر استنتاج مستهلک‌شده معرفی می‌شود. پس از یادگیری سازگاری پرامپت و پاسخ توسط مدل انرژی، چشم‌انداز آن یک شبکه‌ی مفسر سبک را هدایت می‌کند. مفسر با شرطی‌شدن بر روی پرامپت، پاسخ و جملات هدف، تابع هدف شکاف انرژی را بهینه‌سازی می‌کند تا تعادل میان کفایت و ضرورت برقرار شود. در زمان آزمون، این مفسر به صورت مستقل عمل کرده و میزان انتساب جملات پرامپت را در زمان‌ثابت $\mathcal{O}(1)$ و با یک عبور پیش‌رو تخمین می‌زند، بدون آنکه نیازی به استعلام مجدد از مدل اصلی داشته باشد.

ارزیابی‌های تجربی نشان می‌دهد که مدل انرژی با عبور از تطبیق‌های سطحی، الگوهای غیربدیهی تولید را شناسایی می‌کند. همچنین ارزیابی‌های مداخله‌ای مبتنی بر استلزام منطقی، شباهت معنایی و احتمالات توکن‌ها حاکی از آن است که \lr{ESCI} وفاداری معنایی رقابتی در مقایسه با روش‌های پرهزینه‌ی اختلالی و روش‌های ممتاز ارائه می‌دهد. با انتقال بار محاسباتی به فاز پیش‌آموزش، این پژوهش چارچوبی مقیاس‌پذیر برای بررسی مدل‌های مولد جعبه‌سیاه فراهم می‌سازد.

\thesiskeywords{مدل‌های زبانی بزرگ جعبه‌سیاه، تفسیرپذیری پساعملی، انتساب ویژگی در سطح جمله، مدل‌های مبتنی بر انرژی، استنتاج مستهلک‌شده، ممیزی هوش مصنوعی مولد.}
\end{persianabstract}
` in `thesis.tex`.
% ==============================================================================

\begin{persianabstract}
به‌کارگیری مدل‌های زبانی بزرگ از طریق واسط‌های برنامه‌نویسی مسدود در حوزه‌های پرمخاطره، تفسیرپذیری را ضروری ساخته است. با این حال، پارادایم‌های متداول نارسایی‌های اساسی دارند: رویکردهای درونی با تحمیل قیدهای ساختاری، ظرفیت مدل را محدود کرده و برای سیستم‌های منجمد ناکارآمدند؛ روش‌های پساعملی جعبه‌سفید به اطلاعات درونی پنهان در پسِ واسط‌های تجاری وابسته‌اند؛ و روش‌های پساعملی جعبه‌سیاه نیز یا با استدلال‌های هدفمند در خودتوضیحی مواجه‌اند یا در رویکردهای اختلالی توکن‌محور، به فروپاشی دستوری، خطای خارج از توزیع، عدم درک واژگان هم‌معنی و هزینه‌های نمایی دچار می‌شوند. افزون بر این، پیچیدگی ماهوی مدل‌های مولد سبب شده است تا انتساب ویژگی در رژیم کاملاً جعبه‌سیاه و تنها بر مبنای متن خام، همچنان مسئله‌ای بررسی‌نشده در آن‌ها باقی بماند.

برای حل این مسائل، این پایان‌نامه واحد تحلیل را از توکن‌ها به جملات ارتقا می‌دهد تا با حفظ یکپارچگی دستوری در حذف ویژگی‌ها، توضیحاتی خوانا تولید شوند. افزون بر این، یک مدل مبتنی بر انرژی به عنوان جانشین، هماهنگی پرامپت و پاسخ را بدون نیاز به توابع محاسبه‌ناپذیر می‌سنجد. برای آموزش این چشم‌انداز انرژی در غیاب برچسب‌های حقیقت پایه، راهبرد نمونه‌برداری تقابلی دوگانه‌ای طراحی شده است که وفاداری سراسری (تقابل با متون انسانی و مدل‌های خارجی) را با وابستگی‌های محلی پادواقعی (ماسک‌گذاری ساختاریافته‌ی پرامپت و پاسخ) ترکیب می‌کند.

بر این مبنا، چارچوب \lr{ESCI} به عنوان روشی پساعملی و مبتنی بر استنتاج مستهلک‌شده معرفی می‌شود. پس از یادگیری سازگاری پرامپت و پاسخ توسط مدل انرژی، چشم‌انداز آن یک شبکه‌ی مفسر سبک را هدایت می‌کند. مفسر با شرطی‌شدن بر روی پرامپت، پاسخ و جملات هدف، تابع هدف شکاف انرژی را بهینه‌سازی می‌کند تا تعادل میان کفایت و ضرورت برقرار شود. در زمان آزمون، این مفسر به صورت مستقل عمل کرده و میزان انتساب جملات پرامپت را در زمان‌ثابت $\mathcal{O}(1)$ و با یک عبور پیش‌رو تخمین می‌زند، بدون آنکه نیازی به استعلام مجدد از مدل اصلی داشته باشد.

ارزیابی‌های تجربی نشان می‌دهد که مدل انرژی با عبور از تطبیق‌های سطحی، الگوهای غیربدیهی تولید را شناسایی می‌کند. همچنین ارزیابی‌های مداخله‌ای مبتنی بر استلزام منطقی، شباهت معنایی و احتمالات توکن‌ها حاکی از آن است که \lr{ESCI} وفاداری معنایی رقابتی در مقایسه با روش‌های پرهزینه‌ی اختلالی و روش‌های ممتاز ارائه می‌دهد. با انتقال بار محاسباتی به فاز پیش‌آموزش، این پژوهش چارچوبی مقیاس‌پذیر برای بررسی مدل‌های مولد جعبه‌سیاه فراهم می‌سازد.

\thesiskeywords{مدل‌های زبانی بزرگ جعبه‌سیاه، تفسیرپذیری پساعملی، انتساب ویژگی در سطح جمله، مدل‌های مبتنی بر انرژی، استنتاج مستهلک‌شده، ممیزی هوش مصنوعی مولد.}
\end{persianabstract}
` in `thesis.tex`.
% ==============================================================================

\begin{persianabstract}
به‌کارگیری مدل‌های زبانی بزرگ از طریق واسط‌های برنامه‌نویسی مسدود در حوزه‌های پرمخاطره، تفسیرپذیری را ضروری ساخته است. با این حال، پارادایم‌های متداول نارسایی‌های اساسی دارند: رویکردهای درونی با تحمیل قیدهای ساختاری، ظرفیت مدل را محدود کرده و برای سیستم‌های منجمد ناکارآمدند؛ روش‌های پساعملی جعبه‌سفید به اطلاعات درونی پنهان در پسِ واسط‌های تجاری وابسته‌اند؛ و روش‌های پساعملی جعبه‌سیاه نیز یا با استدلال‌های هدفمند در خودتوضیحی مواجه‌اند یا در رویکردهای اختلالی توکن‌محور، به فروپاشی دستوری، خطای خارج از توزیع، عدم درک واژگان هم‌معنی و هزینه‌های نمایی دچار می‌شوند. افزون بر این، پیچیدگی ماهوی مدل‌های مولد سبب شده است تا انتساب ویژگی در رژیم کاملاً جعبه‌سیاه و تنها بر مبنای متن خام، همچنان مسئله‌ای بررسی‌نشده در آن‌ها باقی بماند.

برای حل این مسائل، این پایان‌نامه واحد تحلیل را از توکن‌ها به جملات ارتقا می‌دهد تا با حفظ یکپارچگی دستوری در حذف ویژگی‌ها، توضیحاتی خوانا تولید شوند. افزون بر این، یک مدل مبتنی بر انرژی به عنوان جانشین، هماهنگی پرامپت و پاسخ را بدون نیاز به توابع محاسبه‌ناپذیر می‌سنجد. برای آموزش این چشم‌انداز انرژی در غیاب برچسب‌های حقیقت پایه، راهبرد نمونه‌برداری تقابلی دوگانه‌ای طراحی شده است که وفاداری سراسری (تقابل با متون انسانی و مدل‌های خارجی) را با وابستگی‌های محلی پادواقعی (ماسک‌گذاری ساختاریافته‌ی پرامپت و پاسخ) ترکیب می‌کند.

بر این مبنا، چارچوب \lr{ESCI} به عنوان روشی پساعملی و مبتنی بر استنتاج مستهلک‌شده معرفی می‌شود. پس از یادگیری سازگاری پرامپت و پاسخ توسط مدل انرژی، چشم‌انداز آن یک شبکه‌ی مفسر سبک را هدایت می‌کند. مفسر با شرطی‌شدن بر روی پرامپت، پاسخ و جملات هدف، تابع هدف شکاف انرژی را بهینه‌سازی می‌کند تا تعادل میان کفایت و ضرورت برقرار شود. در زمان آزمون، این مفسر به صورت مستقل عمل کرده و میزان انتساب جملات پرامپت را در زمان‌ثابت $\mathcal{O}(1)$ و با یک عبور پیش‌رو تخمین می‌زند، بدون آنکه نیازی به استعلام مجدد از مدل اصلی داشته باشد.

ارزیابی‌های تجربی نشان می‌دهد که مدل انرژی با عبور از تطبیق‌های سطحی، الگوهای غیربدیهی تولید را شناسایی می‌کند. همچنین ارزیابی‌های مداخله‌ای مبتنی بر استلزام منطقی، شباهت معنایی و احتمالات توکن‌ها حاکی از آن است که \lr{ESCI} وفاداری معنایی رقابتی در مقایسه با روش‌های پرهزینه‌ی اختلالی و روش‌های ممتاز ارائه می‌دهد. با انتقال بار محاسباتی به فاز پیش‌آموزش، این پژوهش چارچوبی مقیاس‌پذیر برای بررسی مدل‌های مولد جعبه‌سیاه فراهم می‌سازد.

\thesiskeywords{مدل‌های زبانی بزرگ جعبه‌سیاه، تفسیرپذیری پساعملی، انتساب ویژگی در سطح جمله، مدل‌های مبتنی بر انرژی، استنتاج مستهلک‌شده، ممیزی هوش مصنوعی مولد.}
\end{persianabstract}
` in `thesis.tex`.
% ==============================================================================

\begin{persianabstract}
به‌کارگیری مدل‌های زبانی بزرگ از طریق واسط‌های برنامه‌نویسی مسدود در حوزه‌های پرمخاطره، تفسیرپذیری را ضروری ساخته است. با این حال، پارادایم‌های متداول نارسایی‌های اساسی دارند: رویکردهای درونی با تحمیل قیدهای ساختاری، ظرفیت مدل را محدود کرده و برای سیستم‌های منجمد ناکارآمدند؛ روش‌های پساعملی جعبه‌سفید به اطلاعات درونی پنهان در پسِ واسط‌های تجاری وابسته‌اند؛ و روش‌های پساعملی جعبه‌سیاه نیز یا با استدلال‌های هدفمند در خودتوضیحی مواجه‌اند یا در رویکردهای اختلالی توکن‌محور، به فروپاشی دستوری، خطای خارج از توزیع، عدم درک واژگان هم‌معنی و هزینه‌های نمایی دچار می‌شوند. افزون بر این، پیچیدگی ماهوی مدل‌های مولد سبب شده است تا انتساب ویژگی در رژیم کاملاً جعبه‌سیاه و تنها بر مبنای متن خام، همچنان مسئله‌ای بررسی‌نشده در آن‌ها باقی بماند.

برای حل این مسائل، این پایان‌نامه واحد تحلیل را از توکن‌ها به جملات ارتقا می‌دهد تا با حفظ یکپارچگی دستوری در حذف ویژگی‌ها، توضیحاتی خوانا تولید شوند. افزون بر این، یک مدل مبتنی بر انرژی به عنوان جانشین، هماهنگی پرامپت و پاسخ را بدون نیاز به توابع محاسبه‌ناپذیر می‌سنجد. برای آموزش این چشم‌انداز انرژی در غیاب برچسب‌های حقیقت پایه، راهبرد نمونه‌برداری تقابلی دوگانه‌ای طراحی شده است که وفاداری سراسری (تقابل با متون انسانی و مدل‌های خارجی) را با وابستگی‌های محلی پادواقعی (ماسک‌گذاری ساختاریافته‌ی پرامپت و پاسخ) ترکیب می‌کند.

بر این مبنا، چارچوب \lr{ESCI} به عنوان روشی پساعملی و مبتنی بر استنتاج مستهلک‌شده معرفی می‌شود. پس از یادگیری سازگاری پرامپت و پاسخ توسط مدل انرژی، چشم‌انداز آن یک شبکه‌ی مفسر سبک را هدایت می‌کند. مفسر با شرطی‌شدن بر روی پرامپت، پاسخ و جملات هدف، تابع هدف شکاف انرژی را بهینه‌سازی می‌کند تا تعادل میان کفایت و ضرورت برقرار شود. در زمان آزمون، این مفسر به صورت مستقل عمل کرده و میزان انتساب جملات پرامپت را در زمان‌ثابت $\mathcal{O}(1)$ و با یک عبور پیش‌رو تخمین می‌زند، بدون آنکه نیازی به استعلام مجدد از مدل اصلی داشته باشد.

ارزیابی‌های تجربی نشان می‌دهد که مدل انرژی با عبور از تطبیق‌های سطحی، الگوهای غیربدیهی تولید را شناسایی می‌کند. همچنین ارزیابی‌های مداخله‌ای مبتنی بر استلزام منطقی، شباهت معنایی و احتمالات توکن‌ها حاکی از آن است که \lr{ESCI} وفاداری معنایی رقابتی در مقایسه با روش‌های پرهزینه‌ی اختلالی و روش‌های ممتاز ارائه می‌دهد. با انتقال بار محاسباتی به فاز پیش‌آموزش، این پژوهش چارچوبی مقیاس‌پذیر برای بررسی مدل‌های مولد جعبه‌سیاه فراهم می‌سازد.

\thesiskeywords{مدل‌های زبانی بزرگ جعبه‌سیاه، تفسیرپذیری پساعملی، انتساب ویژگی در سطح جمله، مدل‌های مبتنی بر انرژی، استنتاج مستهلک‌شده، ممیزی هوش مصنوعی مولد.}
\end{persianabstract}
